% ------------------------------------------------------------------------
% ------------------------------------------------------------------------
% ICMC: Modelo de Trabalho Acadêmico (tese de doutorado, dissertação de
% mestrado e trabalhos monográficos em geral) em conformidade com 
% ABNT NBR 14724:2011: Informação e documentação - Trabalhos acadêmicos -
% Apresentação
% ------------------------------------------------------------------------
% ------------------------------------------------------------------------

% nao modificar
\documentclass[english]{packages/icmc}

% ---
% Pacotes Opcionais
% ---
\usepackage{rotating}           % Usado para rotacionar o texto
\usepackage[all,knot,arc,import,poly]{xy}   % Pacote para desenhos gráficos
% Este pacote pode conflitar com outros pacotes gráficos como o ``pictex''
% Então é necessário usar apenas um dos pacotes conflitantes


% ---
% Informações de dados para CAPA e FOLHA DE ROSTO
% ---
\titulo{Titulo do Trabalho de Conclusão de Curso }
\autor[Sobrenome, N. A.]{Nome do(a) Aluno(a) e Sobrenome}
\orientador[Orientador:]{Prof. Dr.}{Nome do Orietandor(a) e Sobrenome}
%\coorientador{Prof. Dr.}{Fulano de tal}
\curso{Sistemas de Informação}
\area{Engenharia de Software} % Área de concentração do trabalho
\data{05}{10}{2022} % Data do depósito
% ---


% ---
% RESUMOS
% ---

% Resumo em português
% conter no máximo 500 palavras
\textoresumo{
\textcolor{red}{Observação 1: Esta seção é obrigatória.}
\textcolor{red}{Observação 2:  O resumo deve apresentar de forma concisa os pontos relevantes do texto. Deve descrever, de maneira objetiva e sucinta, o propósito da monografia, o objetivo do trabalho, sua relevância para a área de pesquisa onde está inserido, a metodologia empregada, e os resultados obtidos, com isso estimular a consulta ao texto completo da monografia.}
\textcolor{red}{Observação 3:   Deve ser redigido em parágrafo único, através de uma seqüência de frases concisas e objetivas e com espaçamento duplo.}
\textcolor{red}{Observação 4:   O resumo deve ter em torno de 15 linhas.}
\textcolor{red}{Observação 5:  Mais informações sobre a escrita de resumo, consulte a norma NBR 6028/1990 para escrita de resumos.}
    }{Palavra-chave 1, Palavra-chave 2, Palavra-chave 3, Palavra-chave 4, Palavra-chave 5}

% ---
% resumo em inglês
% ---
\textoresumo[english]{
\textcolor{red}{Abstract é opcional}
    }{Keyword 1, Keyword 2, Keyword 3, Keyword 4, Keyword 5}
    
% ---
% Configurações de aparência do PDF final
% ---
% alterando o aspecto da cor azul
\definecolor{blue}{RGB}{41,5,195}

% informações do PDF
\makeatletter
\hypersetup{
        pagebackref=true,
		pdftitle={\@title}, 
		pdfauthor={\@author},
    	pdfsubject={\imprimirpreambulo},
	    pdfcreator={LaTeX with abnTeX2/ICMC-USP},
		pdfkeywords={\palavraschave}, 
		colorlinks=true,       		% false: boxed links; true: colored links
    	linkcolor=blue,          	% color of internal links
    	citecolor=blue,        		% color of links to bibliography
    	filecolor=magenta,      	% color of file links
		urlcolor=blue,
		bookmarksdepth=4
}
\makeatother
% --- 

% ----------------------------------------------------------
% ELEMENTOS PRÉ-TEXTUAIS
% ----------------------------------------------------------

% Inserir a ficha catalográfica
%\incluifichacatalografica[tex/fichaCatalografica.pdf]
\incluifichacatalografica


% Inserir folha de aprovação
%\input{tex/folha-aprovacao}

% DEDICATÓRIA / AGRADECIMENTO / EPÍGRAFE
\textodedicatoria*{tex/pre-textual/dedicatoria}
\textoagradecimentos*{tex/pre-textual/agradecimentos}
\textoepigrafe*{tex/pre-textual/epigrafe}

% Inclui a lista de figuras
\incluilistadefiguras

% Inclui a lista de tabelas
\incluilistadetabelas

% Inclui a lista de quadros
\incluilistadequadros

% Inclui a lista de algoritmos
\incluilistadealgoritmos

% Inclui a lista de códigos
\incluilistadecodigos

% Inclui a lista de siglas e abreviaturas
\incluilistadesiglas

% Inclui a lista de símbolos
\incluilistadesimbolos

% ----
% Início do documento
% ----
\begin{document}

% ----------------------------------------------------------
% ELEMENTOS TEXTUAIS
% ----------------------------------------------------------
\textual

\chapter{Introdução}
\label{chapter:introducao}
\textcolor{red}{INSTRUÇÕES GERAIS:}

\textcolor{red}{Observação 1: Este documento é uma recomendação de formatação da monografia, contendo inclusive itens que são esperados no texto da monografia. Contudo, o aluno tem a liberdade de discutir com seu orientador e inserir outros capítulos/seções ou reestruturá-la, caso seja necessário.}

\textcolor{red}{Observação 2: O texto da monografia (marcado como estilo normal do Word)  deve utilizar fonte Times New Roman tamanho 12pt com espaçamento 1,5 entre linhas, conforme já configurado neste documento. Utilizar papel A4.}

\textcolor{red}{Observação 3: Margens: Superior e esquerda: 3,0 cm; Inferior e direita: 2,0 cm.}

\textcolor{red}{Observação 4: Sugere-se que a monografia tenha em torno de 30  páginas.}

\textcolor{red}{Observação 5: Outras recomendações: (i) Substituir todos os textos marcados em vermelho; e (ii) Seções que forem opcionais e não forem utilizadas na monografia devem ser retiradas.}

\textcolor{red}{Lista de Abreviaturas/Siglas: Esta seção é opcional. Insira-a na monografia caso haja uma quantidade razoável de abreviaturas. O fato de incluir a abreviatura/sigla aqui não o libera da obrigatoriedade de colocar a sigla e seu significado na primeira ocorrência dela no texto. A partir da segunda ocorrência use somente a sigla. Se o leitor se esquecer do significado, pode recorrer a esta lista pra se lembrar, ao invés de ter que ficar procurando no meio do texto.}

\textcolor{red}{Lista de Símbolos: Esta seção é opcional. Insira-a na monografia caso haja uma quantidade razoável de símbolos.}

\textcolor{red}{Lista de Gráficos: Esta seção é opcional. Insira-a na monografia caso haja gráficos ao decorrer do texto.
Sugestão: Utilize o índice de ilustrações (gráficos) do Word para gerar automaticamente esta lista.}

\textcolor{red}{Lista de Tabelas: Esta seção é opcional. Insira-a na monografia caso haja tabelas ao decorrer do texto.
Sugestão: Utilize o índice de ilustrações (tabelas) do Word para gerar automaticamente esta lista.}

\textcolor{red}{Lista de Algoritmos: Este item é opcional. Insira-o na monografia caso haja uma quantidade razoável de algoritmos.}

\textcolor{red}{Lista de Códigos-fonte: Este item é opcional. Insira-o na monografia caso haja uma quantidade razoável de códigos-fonte.}

\textcolor{red}{Lista de Figuras: Esta seção é opcional. Insira-a na monografia caso haja figuras ao decorrer do texto.
Sugestão: Utilize o índice de ilustrações (figuras) do Word para gerar automaticamente esta lista.}

\textcolor{red}{INSTRUÇÕES INTRODUÇÃO:}

\textcolor{red}{Observação: um bom capítulo de introdução tem em torno de 3 a 5 páginas. Sugere-se que a monografia tenha em torno de 30 páginas. Caso tenha muito conteúdo, pode escolher algumas partes e colocar como Apêndice. Substitua todos os textos marcados em vermelho e também as Seções que forem opcionais e não forem utilizadas.}

\section{Motivação e Contextualização}

\textcolor{red}{Nesta seção descrevem-se a área de pesquisa na qual o trabalho está inserido, o problema e/ou as circunstâncias que motivaram o projeto e as potenciais contribuições oriundas de sua realização. Além disso, é necessário sintetizar o que foi feito no Projeto I, caso o aluno esteja matriculado atualmente no Projeto II. Ao longo de todo o texto, se citar figuras, cite-a pelo nome, explique a figura e cite a fonte (se for sua, coloque Fonte: autor desta monografia), como nesse exemplo (veja a seção Corpos Flutuantes Figura \ref{figura:logomarca_usp}). Vejam que a legenda vai acima da Figura e a fonte abaixo. Sempre que se referir à figura, use seu número (e não: figura abaixo, figura acima, figura a seguir, etc). O mesmo vale para Tabela, Gráfico, Algoritmo, etc. Toda figura, tabela, etc. tem que estar citada no texto para chamar a atenção do leitor para ela.}

\section{Objetivos}

\textcolor{red}{Indique claramente quais são os objetivos do trabalho, caracterizando de forma sucinta o que se pretende atingir. Se o emprego de uma metodologia específica de trabalho for relevante, indique também; caso contrário, ela aparece apenas na seção de desenvolvimento do trabalho.}

\section{Organização}

\textcolor{red}{Escreva um parágrafo que resume cada um dos demais capítulos da monografia, fazendo referência a cada um deles – como no exemplo que segue. Indique também a existência de apêndices e anexos, se houver.  No Capítulo 2 é apresentada uma revisão da terminologia básica utilizada no projeto, bem como os trabalhos da literatura relacionados ao presente projeto. A seguir, no Capítulo 3, descrevem-se as atividades realizadas. Finalmente, no Capítulo X, apresentam-se as conclusões e trabalhos futuros. Observe que o tempo verbal é presente, pois refere-se ao texto que já está escrito na monografia. Use tempo futuro somente para se referir a atividades que ainda serão feitas.}

\chapter{Revisão Bibliográfica}
\label{chapter:bibliografia}
\textcolor{red}{(Este capítulo em geral não deve ultrapassar o total de 5 páginas).}

\section{Considerações Iniciais}

\textcolor{red}{Neste capítulo devem ser apresentados os conceitos e a terminologia básicos da área do projeto e o levantamento bibliográfico necessário para a realização do trabalho. Em particular, a descrição dos principais trabalhos de pesquisa relacionados com este (em geral, aqueles que representam o estado da arte na área de pesquisa em questão), bem como dos trabalhos que serviram de base para a solução proposta por este projeto (em geral, aqueles que apresentam técnicas ou recursos que foram utilizados pelo projeto). A divisão nas subseções a seguir é opcional. Em particular, nesta seção (Considerações Iniciais), deve-se descreva sucintamente o que é apresentado neste capítulo.}

\section{Título}
\subsection{Subtítulo 1}
\subsection{Subtítulo 2}
\textcolor{red}{(As seções e subsubseções podem ser organizadas da forma que melhor se adequar à sua monografia)}

\section{Trabalhos Relacionados}

\textcolor{red}{Refereciar aqui os trabalhos relacionados ao projeto. Não esquecer de citar a fonte corretamente, por exemplo [Silva e Zhao 2012].}

\section{Considerações Finais}

\textcolor{red}{Nesta seção deve-se apresentar uma conclusão sobre este capítulo e introduzir brevemente o capítulo seguinte.}

\chapter{Desenvolvimento}
\label{chapter:desenvolvimento}
\input{tex/capitulos/desenvolvimento}

\chapter{Conclusão}
\label{chapter:conclusao}
\textcolor{red}{(Este capítulo em geral não deve ultrapassar o total de 3 páginas.)}

\section{Contribuições}

\textcolor{red}{Nesta seção, o aluno deve destacar as principais contribuições do trabalho realizado, tanto no nível profissional quanto pessoal. Descreva conclusões a respeito do trabalho desenvolvido e experiências adquiridas.}

\section{Considerações Sobre o Curso de Graduação}

\textcolor{red}{Pode começar descrevendo como as disciplinas do curso ou atividades desenvolvidas durante o curso de graduação se relacionam ou contribuíram para o projeto desenvolvido. 
 Depois, coloque aqui as suas críticas ao curso de graduação, pontuando os pontos fracos, (e porque não) os fortes também, que limitaram (ou contribuíram para) suas ações no desenvolvimento do trabalho. Faça sugestões que você achar pertinente para a melhora do curso, disciplinas ou temas novos a serem abordados, etc.  }

\section{Trabalhos Futuros}

\textcolor{red}{Nesta seção, o aluno pode apresentar atividades e sugestões a serem realizadas para a continuação do trabalho realizado neste projeto. Esta seção é opcional.}

\section{Referências (Seção apenas explicativa, não deve ir no trabalho)}

\textcolor{red}{Observação 1: As referências devem estar em ordem alfabética pelo sobrenome do primeiro autor. Observação 2: Todos os documentos referenciados nesta seção (livros, artigos, relatórios técnicos, sites, entre outros) deverão ter sido citados no texto. Observação 3: Todos as citações no decorrer do texto deverão ser listadas nesta seção. Observação 4: Todos as referências devem ser identificáveis quanto a forma de publicação: por exemplo, se for um livro, deve ter a Editora, se for um artigo, deve ter o nome do periódico, se for uma tese ou dissertação, isso dever estar claro e deve ter o nome da Universidade onde foi defendida. Veja a observação 5 a respeito disso, pois seguindo a norma não tem como errar. Observação 5: É obrigatório que a lista de referências esteja de acordo com a norma NBR-6023/2002 para referenciação bibliográfica.}

\chapter{A Partir Deste Capítulo Apenas Orientações Sobre o Uso do Latex}
\label{chapter:orientacoes-gerais}
\input{tex/orientacoes-gerais}

\chapter{Configuração dos Elementos Pré-Textuais}
\label{chapter:config-pre-textual}
\input{tex/config-pre-textual}

\chapter{Corpos Flutuantes}
\label{chapter:corpos-flutuantes}
\input{tex/corpos-flutuantes}

\chapter{Listas}
\label{chapter:listas}
\input{tex/listas}

\chapter{Ferramentas Úteis}
\label{chapter:ferramentas-uteis}
\input{tex/ferramentas-uteis}

\chapter{Citações e Referências}
\label{chapter:citacoes}
\input{tex/citacoes}


% ---
% Finaliza a parte no bookmark do PDF, para que se inicie o bookmark na raiz
% ---
\bookmarksetup{startatroot}% 
% ---

% ----------------------------------------------------------
% ELEMENTOS PÓS-TEXTUAIS
% ----------------------------------------------------------
\postextual

% ----------------------------------------------------------
% Referências bibliográficas
% ----------------------------------------------------------
\bibliography{references}

% ---------------------------------------------------------------------
% GLOSSÁRIO
% ---------------------------------------------------------------------

% Arquivo que contém as definições que vão aparecer no glossário
\input{tex/glossario}
% Comando para incluir todas as definições do arquivo glossario.tex
\glsaddall
% Impressão do glossário
\printglossaries

% ----------------------------------------------------------
% Apêndices
% ----------------------------------------------------------

% ---
% Inicia os apêndices
% ---
\begin{apendicesenv}

    \chapter{Título do Apêndice}
    \label{chapter:tituloapendice}
    \textcolor{red}{Observação 1: Apêndice consiste em um texto ou documento elaborado pelo autor, a fim de complementar sua argumentação, sem prejuízo da unidade nuclear do trabalho. Os apêndices são identificados por letras maiúsculas consecutivas, travessão e pelos respectivos títulos. Observação 2: Elemento opcional. Observação 3: Se houverem mais apêndices, identifique-os como Apêndice B, Apêndice C e assim por diante.}

\end{apendicesenv}
% ---


% ----------------------------------------------------------
% Anexos
% ----------------------------------------------------------

% ---
% Inicia os anexos
% ---
\begin{anexosenv}

    \chapter{Título do Anexo} 
    \label{chapter:tituloanexo}
    \input{tex/annex/exemploanexo}

\end{anexosenv}
% ---

\end{document}